\documentclass[slidestop,compress,mathserif,c]{beamer}
\usepackage{listings}
\usepackage{color}
\usepackage{xcolor}
\usepackage{hyperref}
\definecolor{dkgreen}{rgb}{0,0.6,0}
\usepackage{fontspec,xunicode,xltxtra,beamerthemesplit}
%\usepackage{slashbox}
\usepackage{diagbox}
 \newfontfamily\times{Times New Roman}
%以下是各种演示主题，定义幻灯片中的所有细节
%\usetheme{default}
%\usetheme{Berlin}%这个主题比较好
%\usetheme{Pittsburgh}
%\usetheme{Rochester}
\usetheme{Berkeley}%这种演示主题比较好
%\usetheme{Goettingen}
%\usetheme{Hannover}
%\usetheme{Marburg}
%\usetheme{PaloAlto}%这种演示主题比较好
%\usetheme{Antibes}
%\usetheme{Darmstadt}
%\usetheme{JuanLesPins}
%\usetheme{Montpellier}
%\usetheme{Singapore}
%\usetheme{Boadilla}
%\usetheme{Madrid}
%\usetheme{AnnArbor}
%\usetheme{CambridgeUS}
%\usetheme{Copenhagen}
%\usetheme{Warsaw}

%以下是各种外部主题，确定幻灯的显示式样
%\useoutertheme{infolines}
%\useoutertheme{miniframes}
%\useoutertheme[height=0.1\textwidth,width=0.15\textwidth,hideothersubsections]{sidebar}
%\useoutertheme{smoothbars}
%\useoutertheme{split}
%\useoutertheme{shadow}
%\useoutertheme{tree}
%\useoutertheme{smoothtree}
%\useoutertheme[height=0.5\textwidth]{sidebar}

%以下是各种内部主题
%\useinnertheme{default}
%\useinnertheme{circles}
%\useinnertheme{rectangles}
\useinnertheme[shadow]{rounded}

%以下是各种颜色主题
%\usecolortheme{default}
%\usecolortheme{albatross}
%\usecolortheme{beaver}
%\usecolortheme{beetle}
%\usecolortheme{crane}
%\usecolortheme{dolphin}
%\usecolortheme{dove}
%\usecolortheme{fly}
%\usecolortheme{lily}
%\usecolortheme{orchid}
%\usecolortheme{rose}
%\usecolortheme{seagull}
\usecolortheme{seahorse}
%\usecolortheme{sidebartab}
%\usecolortheme{structure}
%\usecolortheme{whale}
%\usecolortheme{wolverine}

%以下是各种字体主题
%\usefonttheme{default}
%\usefonttheme[onlymath]{serif}
%\usefonttheme{structurebold}
%\usefonttheme{structureitalicserif}
%\usefonttheme{structuresmallcapsserif}
\newcommand{\tablecell}[2]{\begin{tabular}{@{}#1@{}}#2\end{tabular}}   % 表格内换行，用法\tablecell{c}{balabalba\\balabalba}
\setsansfont[Mapping=tex-text, BoldFont={微软雅黑 Bold}]{微软雅黑}

% 中文环境自动换行
\XeTeXlinebreaklocale "zh"  % 表示用中文的断行
\XeTeXlinebreakskip = 0pt plus 1pt % 多一点调整的空间

% 中文环境修正导航栏
\makeatletter
\setbeamertemplate{blocks}[rounded][shadow=true] 
\def\beamer@linkspace#1{
  \begin{pgfpicture}{0pt}{-1.5pt}{#1}{5.5pt}
    \pgfsetfillopacity{0}
    \pgftext[x=0pt,y=-1.5pt]{.}
    \pgftext[x=#1,y=5.5pt]{.}
  \end{pgfpicture}}
\makeatother

% 超链接高亮显示
\hypersetup{CJKbookmarks=true,
colorlinks=true,
citecolor=blue,
linkcolor=blue,
urlcolor=blue,
bookmarksopen=true,
breaklinks=true
}

% 幻灯片切换方式
%\transblindshorizontal 	% 水平百叶窗
%\transblindsvertical 		% 垂直百叶窗
%\transboxin				% 盒状收缩
%\transboxout				% 盒状展开
%\transdissolve				% 溶解
%\transglitter
%\transsplithorizontalin	% 上下向中央收缩
%\transsplitverticalin		% 垂直向中央收缩
%\transsplithorizontalout	% 上下向中央展开
%\transsplitverticalout		% 垂直向中央展开
%\transwipe					% 从下抽出


\title{Gossip算法}
\author{Liu Zheng}
\date{\today}
\institute{同济大学电信学院}

\logo{\color{blue!50}\scalebox{2}{{\includegraphics[height=0.8cm]{logo.jpg}\vspace{220pt}}}}

% 设定frametitle居中
\makeatletter 
\long\def\beamer@@frametitle[#1]#2{% 
  \beamer@ifempty{#2}{}{% 
    \gdef\insertframetitle{\centering{#2\ifnum\beamer@autobreakcount>0\relax{}\space\usebeamertemplate*{frametitle continuation}\fi}}% 
  \gdef\beamer@frametitle{#2}% 
  \gdef\beamer@shortframetitle{#1}% 
}% 
} 
\makeatother


\begin{document}

\frame{\titlepage}

\section{Gossip背景}

%%%%%%%%%% One Slide %%%%%%%%%%
%\subsection{\hfill Overview}
\begin{frame}
\frametitle{Overview}
~~~~~~Gossip算法因为Cassandra而名声大噪，Gossip看似简单，但要真正弄清楚其本质远没看起来那么容易。为了寻求Gossip的本质，下面的内容主要参考Gossip的原始论文：《Efficient Reconciliation and Flow Control for Anti-Entropy Protocols》。
\end{frame}
%%%%%%%%%% Slide End %%%%%%%%%%

%%%%%%%%%% One Slide %%%%%%%%%%
\subsection{\hfill  Gossip背景}
\begin{frame}
\frametitle{Gossip背景}
~~~~~~Gossip算法如其名，灵感来自办公室八卦，只要一个人八卦一下，在有限的时间内所有的人都会知道该八卦的信息，这种方式也与病毒传播类似，因此Gossip有众多的别名“闲话算法”、“疫情传播算法”、“病毒感染算法”、“谣言传播算法”。

~~~~~~但Gossip并不是一个新东西，之前的泛洪查找、路由算法都归属于这个范畴，不同的是Gossip给这类算法提供了明确的语义、具体实施方法及收敛性证明。
\end{frame}
%%%%%%%%%% Slide End %%%%%%%%%%


%%%%%%%%%% One Slide %%%%%%%%%%
\subsection{\hfill Gossip算法的发展历史}
\begin{frame}
\frametitle{Gossip算法的发展历史}
 
\end{frame}
%%%%%%%%%% Slide End %%%%%%%%%%


\section{算法描述}

%%%%%%%%%% One Slide %%%%%%%%%%
%\subsection{\hfill hehe}
\begin{frame}
\frametitle{算法描述}
~~~~~~假设有 ${p, q, \cdots}$ 为协议参与者。 每个参与者都有关于一个自己信息的表。
用编程语言可以描述为： 
记 InfoMap = Map<Key, (Value, Version)>， 那么每个参与者要维护一个 InfoMap 类型的变量 localInfo。 同时每一个参与者要知道所有其他参与者的信息, 即要维护一个全局的表，即 Map<participant, InfoMap> 类型的变量 globalMap。每个参与者更新自己的 localInfo， 而由 Gossip 协议负责将更新的信息同步到整个网络上。
每个节点和系统中的某些节点成为 peer (如果系统的规模比较小，和系统中所有的其他节点成为 peer)。 
\end{frame}
%%%%%%%%%% Slide End %%%%%%%%%%


%%%%%%%%%% One Slide %%%%%%%%%%
%\subsection{\hfill hehe}
\begin{frame}
\frametitle{hehe}
 三种不同的同步信息的方法：
 \begin{itemize}
\item[1）]push-gossip: 最简单的情况下， 一个节点 p 向 q 发送整个 GlobalMap
\item[2）]pull-gossip: p 向 q 发送 digest, q 根据 digest 向 p 发送 p 过期的 (key, (value, version)) 列表
\item[3）]push-pull-gossip:与pull-gossip类似，只是多了一步，A再将本地比B新的数据推送给B，B更新本地 
 \end{itemize}
\end{frame}
%%%%%%%%%% Slide End %%%%%%%%%%


%%%%%%%%%% One Slide %%%%%%%%%%
\subsection{\hfill Gossip特点}
\begin{frame}
\frametitle{Gossip特点}
~~~~~~Gossip算法又被称为反熵（Anti-Entropy），熵是物理学上的一个概念，代表杂乱无章，而反熵就是在杂乱无章中寻求一致，这充分说明了Gossip的特点：在一个有界网络中，每个节点都随机地与其他节点通信，经过一番杂乱无章的通信，最终所有节点的状态都会达成一致。每个节点可能知道所有其他节点，也可能仅知道几个邻居节点，只要这些节可以通过网络连通，最终他们的状态都是一致的，当然这也是疫情传播的特点。

~~~~~~要注意到的一点是，即使有的节点因宕机而重启，有新节点加入，但经过一段时间后，这些节点的状态也会与其他节点达成一致，也就是说，Gossip天然具有分布式容错的优点。
\end{frame}
%%%%%%%%%% Slide End %%%%%%%%%%

%%%%%%%%%% One Slide %%%%%%%%%%
\subsection{\hfill Gossip本质}
\begin{frame}
\frametitle{Gossip本质}
~~~~~~Gossip是一个带冗余的容错算法，更进一步，Gossip是一个最终一致性算法。虽然无法保证在某个时刻所有节点状态一致，但可以保证在”最终“所有节点一致，”最终“是一个现实中存在，但理论上无法证明的时间点。

~~~~~~因为Gossip不要求节点知道所有其他节点，因此又具有去中心化的特点，节点之间完全对等，不需要任何的中心节点。实际上Gossip可以用于众多能接受“最终一致性”的领域：失败检测、路由同步、Pub/Sub、动态负载均衡。

~~~~~~但Gossip的缺点也很明显，冗余通信会对网路带宽、CUP资源造成很大的负载，而这些负载又受限于通信频率，该频率又影响着算法收敛的速度，后面我们会讲在各种场合下的优化方法。
\end{frame}
%%%%%%%%%% Slide End %%%%%%%%%%

%%%%%%%%%% One Slide %%%%%%%%%%
\subsection{\hfill  Gossip节点的通信方式及收敛性}
\begin{frame}
\frametitle{ Gossip节点的通信方式及收敛性}
 根据原论文，两个节点 A、B 之间存在三种通信方式:
 \begin{itemize}
 \item push: A节点将数据(key,value,version)及对应的版本号推送给B节点，B节点更新A中比自己新的数据
 \item pull：A仅将数据key,version推送给B，B将本地比A新的数据（Key,value,version）推送给A，A更新本地
 \item push/pull：与pull类似，只是多了一步，A再将本地比B新的数据推送给B，B更新本地
 \end{itemize}
\end{frame}
%%%%%%%%%% Slide End %%%%%%%%%%

%%%%%%%%%% One Slide %%%%%%%%%%
%\subsection{\hfill hehe}
\begin{frame}
\frametitle{hehe}
 ~~~~~~如果把两个节点数据同步一次定义为一个周期，则在一个周期内，push需通信1次，pull需2次，push/pull则需3次，从效果上来讲，push/pull最好，理论上一个周期内可以使两个节点完全一致。直观上也感觉，push/pull的收敛速度是最快的。

~~~~~~假设每个节点通信周期都能选择（感染）一个新节点，则Gossip算法退化为一个二分查找过程，每个周期构成一个平衡二叉树，收敛速度为$O(n^2)$，对应的时间开销则为$O(\log n)$。这也是Gossip理论上最优的收敛速度。但在实际情况中最优收敛速度是很难达到的，假设某个节点在第i个周期被感染的概率为$p_i$ ,第$i+1$个周期被感染的概率为$p_{i+1}$ ，则pull的方式:
~~~~~~~~~~~~~~~~$ p_{i+1} =p_i^2$

而push为：~~~~~~~~~~~~
$p_{i+1}=p_i(1-\frac{1}{n})^{n(1-p_i)}$
\end{frame}
%%%%%%%%%% Slide End %%%%%%%%%%

%%%%%%%%%% One Slide %%%%%%%%%%
%\subsection{\hfill hehe}
\begin{frame}
\frametitle{hehe}
 ~~~~~~显然pull的收敛速度大于push，而每个节点在每个周期被感染的概率都是固定的$p(0<p<1)$，因此Gossip算法是基于p的平方收敛，也成为概率收敛，这在众多的一致性算法中是非常独特的。
\end{frame}
%%%%%%%%%% Slide End %%%%%%%%%%

%%%%%%%%%% One Slide %%%%%%%%%%
%\subsection{\hfill hehe}
\begin{frame}
\frametitle{hehe}
 Gossip的节点的工作方式又分两种：
 \begin{itemize}
 \item Anti-Entropy（反熵）：以固定的概率传播所有的数据
 \item Rumor-Mongering（谣言传播）：仅传播新到达的数据
 \end{itemize}
 ~~~~~~Anti-Entropy模式有完全的容错性，但有较大的网络、CPU负载；Rumor-Mongering模式有较小的网络、CPU负载，但必须为数据定义”最新“的边界，并且难以保证完全容错，对失败重启且超过”最新“期限的节点，无法保证最终一致性，或需要引入额外的机制处理不一致性。我们后续着重讨论Anti-Entropy模式的优化。
\end{frame}
%%%%%%%%%% Slide End %%%%%%%%%%

%%%%%%%%%% One Slide %%%%%%%%%%
\subsection{\hfill Anti-Entropy的协调机制}
\begin{frame}
\frametitle{Anti-Entropy的协调机制}
~~~~~~ 协调机制是讨论在每次2个节点通信时，如何交换数据能达到最快的一致性，也即消除两个节点的不一致性。上面所讲的push、pull等是通信方式，协调是在通信方式下的数据交换机制。协调所面临的最大问题是，因为受限于网络负载，不可能每次都把一个节点上的数据发送给另外一个节点，也即每个Gossip的消息大小都有上限。在有限的空间上有效率地交换所有的消息是协调要解决的主要问题。
\end{frame}
%%%%%%%%%% Slide End %%%%%%%%%%

%%%%%%%%%% One Slide %%%%%%%%%%
%\subsection{\hfill hehe}
\begin{frame}
\frametitle{hehe}
 在讨论之前先声明几个概念：
 \begin{itemize}
 \item 令N = {p,q,s,...}为需要gossip通信的server集合，有界大小
 \item 令(p1,p2,...)是宿主在节点p上的数据，其中数据有(key,value,version)构成，q的规则与p类似。
 \end{itemize}
~~~~~~ 为了保证一致性，规定数据的value及version只有宿主节点才能修改，其他节点只能间接通过Gossip协议来请求数据对应的宿主节点修改。
\end{frame}
%%%%%%%%%% Slide End %%%%%%%%%%

%%%%%%%%%% One Slide %%%%%%%%%%
%\subsection{\hfill hehe}
\begin{frame}
\frametitle{精确协调（Precise Reconciliation）}
 ~~~~~~精确协调希望在每次通信周期内都非常准确地消除双方的不一致性，具体表现为相互发送对方需要更新的数据，因为每个节点都在并发与多个节点通信，理论上精确协调很难做到。精确协调需要给每个数据项独立地维护自己的version，在每次交互是把所有的(key,value,version)发送到目标进行比对，从而找出双方不同之处从而更新。但因为Gossip消息存在大小限制，因此每次选择发送哪些数据就成了问题。当然可以随机选择一部分数据，也可确定性的选择数据。对确定性的选择而言，可以有最老优先（根据版本）和最新优先两种，最老优先会优先更新版本最新的数据，而最新更新正好相反，这样会造成老数据始终得不到机会更新，也即饥饿。
当然，开发这也可根据业务场景构造自己的选择算法，但始终都无法避免消息量过多的问题。
\end{frame}
%%%%%%%%%% Slide End %%%%%%%%%%

%%%%%%%%%% One Slide %%%%%%%%%%
%\subsection{\hfill hehe}
\begin{frame}
\frametitle{整体协调（Scuttlebutt Reconciliation）}
~~~~~~ 整体协调与精确协调不同之处是，整体协调不是为每个数据都维护单独的版本号，而是为每个节点上的宿主数据维护统一的version。比如节点P会为(p1,p2,...)维护一个一致的全局version，相当于把所有的宿主数据看作一个整体，当与其他节点进行比较时，只需必须这些宿主数据的最高version，如果最高version相同说明这部分数据全部一致，否则再进行精确协调。
整体协调对数据的选择也有两种方法：
\begin{itemize}
\item 广度优先：根据整体version大小排序，也称为公平选择
\item 深度优先：根据包含数据多少的排序，也称为非公平选择。因为后者更有实用价值，所以原论文更鼓励后者
\end{itemize}
\end{frame}
%%%%%%%%%% Slide End %%%%%%%%%%

%%%%%%%%%% One Slide %%%%%%%%%%
\subsection{\hfill Cassandra中的实现}
\begin{frame}
\frametitle{Cassandra中的实现}
 经过验证，Cassandra实现了基于整体协调的push/push模式，有几个组件：
三条消息分别对应push/pull的三个阶段：
\begin{itemize}
\item GossipDigitsMessage
\item GossipDigitsAckMessage
\item GossipDigitsAck2Message
\end{itemize}
\end{frame}
%%%%%%%%%% Slide End %%%%%%%%%%

%%%%%%%%%% One Slide %%%%%%%%%%
%\subsection{\hfill hehe}
\begin{frame}
\frametitle{hehe}
 还有三种状态：
 \begin{itemize}
\item EndpointState：维护宿主数据的全局version，并封装了HeartBeat和\item ApplicationState
\item HeartBeat：心跳信息
\item ApplicationState：系统负载信息（磁盘使用率）
 \end{itemize}
\end{frame}
%%%%%%%%%% Slide End %%%%%%%%%%

%%%%%%%%%% One Slide %%%%%%%%%%
%\subsection{\hfill hehe}
\begin{frame}
\frametitle{hehe}
 Cassandra主要是使用Gossip完成三方面的功能：
\begin{itemize}
\item 失败检测
\item 动态负载均衡
\item 去中心化的弹性扩展
\end{itemize}
\end{frame}
%%%%%%%%%% Slide End %%%%%%%%%%

%%%%%%%%%% One Slide %%%%%%%%%%
\subsection{\hfill 总结}
\begin{frame}
\frametitle{总结}
~~~~~~ Gossip是一种去中心化、容错而又最终一致性的绝妙算法，其收敛性不但得到证明还具有指数级的收敛速度。使用Gossip的系统可以很容易的把Server扩展到更多的节点，满足弹性扩展轻而易举。
 
唯一的缺点是收敛是最终一致性，不使用那些强一致性的场景，比如2pc。
\end{frame}
%%%%%%%%%% Slide End %%%%%%%%%%

%%%%%%%%%% One Slide %%%%%%%%%%
%\subsection{\hfill hehe}
\begin{frame}
\frametitle{hehe}
 
\end{frame}
%%%%%%%%%% Slide End %%%%%%%%%%



%%%%%%%%%% One Slide %%%%%%%%%%
%\subsection{\hfill hehe}
\begin{frame}
\frametitle{hehe}
 
\end{frame}
%%%%%%%%%% Slide End %%%%%%%%%%


%%%%%%%%%% One Slide %%%%%%%%%%
%\subsection{\hfill hehe}
\begin{frame}
\frametitle{hehe}
 
\end{frame}
%%%%%%%%%% Slide End %%%%%%%%%%


%%%%%%%%%% One Slide %%%%%%%%%%
%\subsection{\hfill hehe}
\begin{frame}
\frametitle{hehe}
 
\end{frame}
%%%%%%%%%% Slide End %%%%%%%%%%




\end{document}
